%==============================================================================
\section{Hồi quy tuyến tính đơn (Simple Linear Regression)}
\label{sec:simple-linear-regression}
%==============================================================================

\begin{dinhnghia}[title={Hồi quy tuyến tính đơn}]
Phương pháp thống kê và học có giám sát: dùng \textbf{một} biến độc lập (predictor) để dự đoán biến phụ thuộc; mô hình hóa quan hệ tuyến tính giữa hai biến.
\end{dinhnghia}

\begin{phuongphap}[title={Liên hệ với hồi quy tuyến tính}]
Hồi quy tuyến tính đơn là một loại hồi quy tuyến tính.
Thuật toán hồi quy tuyến tính với \textbf{một} feature $\Rightarrow$ hồi quy đơn; với \textbf{nhiều} feature $\Rightarrow$ hồi quy đa biến.
\end{phuongphap}

\subsection{Biến độc lập và phụ thuộc}

\begin{dinhnghia}[title={Biến độc lập (Independent Variable)}]
Feature trong dataset; trong hồi quy đơn chỉ có \textbf{một} biến độc lập.
Còn gọi predictor --- dùng để dự đoán target; vẽ trên trục \textbf{ngang}.
\end{dinhnghia}

\begin{dinhnghia}[title={Biến phụ thuộc (Dependent Variable)}]
Giá trị target trong dataset; còn gọi response hoặc predicted variable; vẽ trên trục \textbf{dọc}.
\end{dinhnghia}

\begin{dinhnghia}[title={Đường hồi quy (Line of Regression)}]
Trong hồi quy đơn, đường hồi quy là đường thẳng khớp tốt nhất các điểm, thể hiện quan hệ giữa biến phụ thuộc và độc lập.
\end{dinhnghia}

\subsection{Biểu diễn đồ thị}

\tsimg{03_simple_linear_regression/img01.jpg}

\subsection{Mô hình toán học}

\begin{dinhnghia}[title={Phương trình}]
\[
Y = w_0 + w_1 X + \epsilon
\]
\begin{itemize}
  \item $Y$: biến phụ thuộc (target);
  \item $X$: biến độc lập (feature);
  \item $w_0$: hệ số chặn;
  \item $w_1$: độ dốc, ảnh hưởng của $X$ lên $Y$;
  \item $\epsilon$: sai số.
\end{itemize}
\end{dinhnghia}

\subsection{Cách hồi quy tuyến tính đơn hoạt động}

\begin{dongco}[title={Mục tiêu}]
Tìm đường thẳng khớp tốt nhất qua các điểm sao cho chênh lệch giữa giá trị thực và dự đoán là nhỏ nhất.
\end{dongco}

\begin{dinhnghia}[title={Hàm giả thuyết}]
Giả thuyết: quan hệ tuyến tính giữa target (đầu ra) và feature (đầu vào):
\[
\hat{Y} = w_0 + w_1 X
\]
Với mỗi cặp $(w_0, w_1)$ có một đường thẳng; tập mọi đường thẳng đó gọi là \textbf{không gian giả thuyết} (hypothesis space).
Mục tiêu: tìm đường khớp tốt nhất trong không gian giả thuyết.
\end{dinhnghia}

\begin{phuongphap}[title={Tìm đường khớp tốt nhất}]
Định nghĩa hàm cost/loss đo chênh lệch giữa giá trị thực và dự đoán.
Mô hình khởi tạo tham số (mặc định), dùng đường hồi quy ban đầu để dự đoán trên đầu vào.
\end{phuongphap}

\begin{dinhnghia}[title={Hàm mất mát (MSE)}]
Dùng giá trị đầu vào và dự đoán để tính loss; loss giúp tìm tham số tối ưu.
Có thể dùng MSE, MAE, $R^2$, v.v.; phổ biến nhất là \textbf{mean squared error}:
\[
J(w_0, w_1) = \frac{1}{2n} \sum_{i=1}^{n} \left( Y_i - \hat{Y}_i \right)^2
\]
\end{dinhnghia}

\begin{phuongphap}[title={Tối ưu hóa}]
Tham số tối ưu là bộ giá trị \textbf{tối thiểu hóa} cost; quá trình lặp cập nhật tham số.
Gradient Descent là kỹ thuật tối ưu đơn giản và phổ biến nhất trong hồi quy đơn.
Phương trình tuyến tính với tham số tối ưu là đường hồi quy cuối cùng, dùng dự đoán dữ liệu mới.
\end{phuongphap}

\subsection{Giả định}

\begin{luuy}[title={Giả định mô hình hồi quy đơn}]
\begin{itemize}
  \item \textbf{Tuyến tính (Linearity)}: biến phụ thuộc thay đổi tuyến tính theo biến độc lập; scatter plot gần đường thẳng.
  \item \textbf{Phương sai đồng nhất (Homoskedasticity)}: phương sai phần dư gần như nhau cho mọi quan sát.
  \item \textbf{Độc lập (Independence)}: các cặp $(X,Y)$ độc lập; kiểm tra bằng scatter phần dư so với fitted values.
  \item \textbf{Chuẩn phần dư (Normality)}: phần dư (thực $-$ dự đoán) gần phân phối chuẩn; kiểm tra histogram phần dư.
\end{itemize}
\end{luuy}

\subsection{Triển khai bằng Python}

\begin{dongco}[title={Dataset và mục đích}]
Hai biến: \texttt{YearsExperience} (độc lập) và \texttt{Salary} (phụ thuộc).
File \path{data/Salary_Data.csv}, 30 mẫu (bảng đầy đủ trong file CSV).
Mục đích: xác định \textbf{đường thẳng nào} biểu diễn tốt nhất quan hệ giữa hai biến.
\end{dongco}

\begin{phuongphap}[title={Bảng dữ liệu (trích)}]
\begin{center}
\small
\begin{tabular}{cc}
\textbf{Năm KN} & \textbf{Lương} \\
1.1 & 39343 \\
1.3 & 46205 \\
\ldots & \ldots \\
10.5 & 121872 \\
\end{tabular}
\end{center}
\end{phuongphap}

\begin{vidu}[title={Bước 1: Chuẩn bị dữ liệu (Data Preparation)}]
Tiền xử lý là bước đầu. Import thư viện trước khi nạp dữ liệu:

\begin{lstlisting}
import numpy as np
import matplotlib.pyplot as plt
import pandas as pd

dataset = pd.read_csv('data/Salary_Data.csv')
\end{lstlisting}

Tách feature và target (\texttt{YearsExperience} = độc lập, \texttt{Salary} = phụ thuộc; bài gốc đảo nhầm X/Y):

\begin{lstlisting}
X = dataset.iloc[:, :-1].values
y = dataset.iloc[:, -1].values
print(dataset.head())
\end{lstlisting}

\textbf{Output mẫu \texttt{head()}:}
\begin{verbatim}
   YearsExperience    Salary
0              1.1   39343.0
1              1.3   46205.0
2              1.5   37731.0
3              2.0   43525.0
4              2.2   39891.0
\end{verbatim}
\end{vidu}

\begin{vidu}[title={Kiểm tra tuyến tính của dataset}]
\begin{lstlisting}
plt.scatter(X, y, color="green")
plt.title("Salary vs Experience")
plt.xlabel("Years of Experience")
plt.ylabel("Salary (INR)")
plt.show()
\end{lstlisting}

\tsimg{03_simple_linear_regression/img02.jpg}

\begin{ynghia}[title={Kết luận từ biểu đồ}]
Biến phụ thuộc và độc lập có quan hệ gần tuyến tính $\Rightarrow$ có thể áp dụng hồi quy tuyến tính đơn.
\end{ynghia}
\end{vidu}

\begin{vidu}[title={Chia tập huấn luyện và kiểm tra}]
30 quan sát: 80\% huấn luyện (24), 20\% kiểm tra (6); một tập train, một tập test.

\begin{lstlisting}
from sklearn.model_selection import train_test_split
X_train, X_test, y_train, y_test = train_test_split(
    X, y, test_size=0.2, random_state=0)
\end{lstlisting}

\texttt{X\_train}, \texttt{y\_train}: feature và target tập huấn luyện.
\end{vidu}

\begin{vidu}[title={Bước 2: Huấn luyện mô hình}]
Dùng \texttt{LinearRegression} của scikit-learn:

\begin{lstlisting}
from sklearn.linear_model import LinearRegression
regressor = LinearRegression()
regressor.fit(X_train, y_train)
\end{lstlisting}

\texttt{fit()} huấn luyện regressor: học quan hệ giữa \texttt{X\_train} và \texttt{y\_train}.
\end{vidu}

\begin{vidu}[title={Bước 3: Kiểm thử mô hình}]
Dự đoán trên tập kiểm tra:

\begin{lstlisting}
y_pred = regressor.predict(X_test)
df = pd.DataFrame({'Actual Values': y_test,
                   'Predicted Values': y_pred})
print(df)
\end{lstlisting}

\texttt{X\_test}: feature kiểm tra; \texttt{y\_pred}: target dự đoán.
Có thể dự đoán tương tự trên tập huấn luyện bằng \texttt{predict(X\_train)}.
\end{vidu}

\begin{vidu}[title={Bước 4: Đánh giá mô hình}]
Dùng MSE, RMSE, MAE, $R^2$:

\begin{lstlisting}
from sklearn.metrics import mean_squared_error, mean_absolute_error, r2_score
import numpy as np

y_pred = regressor.predict(X_test)
print("MSE:", mean_squared_error(y_test, y_pred))
print("RMSE:", np.sqrt(mean_squared_error(y_test, y_pred)))
print("MAE:", mean_absolute_error(y_test, y_pred))
print("R2:", r2_score(y_test, y_pred))
\end{lstlisting}
\end{vidu}

\begin{vidu}[title={Bước 5: Trực quan tập huấn luyện}]
Scatter giá trị thực (xanh) và dự đoán (xanh dương); đường đỏ là đường hồi quy:

\begin{lstlisting}
y_pred_train = regressor.predict(X_train)
plt.scatter(X_train, y_train, color="green",
            label="training data points (actual)")
plt.scatter(X_train, y_pred_train, color="blue",
            label="training data points (predicted)")
plt.plot(X_train, y_pred_train, color="red")
plt.title("Salary vs Experience (Training Dataset)")
plt.xlabel("Years of Experience")
plt.ylabel("Salary (INR)")
plt.legend()
plt.show()
\end{lstlisting}

\tsimg{03_simple_linear_regression/img03.jpg}
\end{vidu}

\begin{vidu}[title={Bước 6: Trực quan tập kiểm tra}]
\begin{lstlisting}
y_pred = regressor.predict(X_test)
plt.scatter(X_test, y_test, color="green",
            label="test data points (actual)")
plt.scatter(X_test, y_pred, color="blue",
            label="test data points (predicted)")
plt.plot(X_test, y_pred, color="red")
plt.title("Salary vs Experience (Test Dataset)")
plt.xlabel("Years of Experience")
plt.ylabel("Salary (INR)")
plt.legend()
plt.show()
\end{lstlisting}

\tsimg{03_simple_linear_regression/img04.jpg}
\end{vidu}
